\paragraph{Pasul 1: V=e, W=t.} V și W sunt cele mai frecvente litere, ca E și T în engleză.
\begin{lstlisting}[language={}]
CIXeOZTG OXO PNZe JeGetXHP RNIL CNI QXZ, ADt, AeHTDPe Qe RTP QTGOFRXtQ T HTZeIT, Qe QeSU
HIFUtNSNJXPtP RQN ReIe SNNLXGJ CNI ATHNG'PHXUQeI PXJGTtDIeP XG PQTLePUeTIe AF ZTLXGJ UQN
eGSTIJeZeGtPNC tQe eSXMTAetQTG UIXGtXGJ tQTt CXJDIeO XG tQe RNIL. tQe OeUTItZeGt
NCHXUQeIP NC tQe IXKeIATGL STANITtNIXeP HNGPXPteO NC 14 NI 15 QXJQ-PHQNNSTGO HNSSeJe JIT
RQN TPPXJGeO tQe XGOXKXODTS SetteIP XG tQePeeSXMTAetQTG teYtP tN NGe NI tQe NtQeI NC tRN
\end{lstlisting}

\paragraph{Pasul 2: Q=h.} trigrama cea mai frecventă, WQV, este aproape sigur THE, deci Q = h.
\begin{lstlisting}[language={}]
CIXeOZTG OXO PNZe JeGetXHP RNIL CNI hXZ, ADt, AeHTDPe he RTP hTGOFRXth T HTZeIT, he heSU
HIFUtNSNJXPtP RhN ReIe SNNLXGJ CNI ATHNG'PHXUheI PXJGTtDIeP XG PhTLePUeTIe AF ZTLXGJ UhN
eGSTIJeZeGtPNC the eSXMTAethTG UIXGtXGJ thTt CXJDIeO XG the RNIL. the OeUTItZeGt
NCHXUheIP NC the IXKeIATGL STANITtNIXeP HNGPXPteO NC 14 NI 15 hXJh-PHhNNSTGO HNSSeJe JIT
RhN TPPXJGeO the XGOXKXODTS SetteIP XG thePeeSXMTAethTG teYtP tN NGe NI the NtheI NC tRN
\end{lstlisting}

\paragraph{Pasul 3: T=a, G=n, O=d.} TGO (trigramă frecventă, cuvânt de 3 litere) este AND.
\begin{lstlisting}[language={}]
CIXedZan dXd PNZe JenetXHP RNIL CNI hXZ, ADt, AeHaDPe he RaP handFRXth a HaZeIa, he heSU
HIFUtNSNJXPtP RhN ReIe SNNLXnJ CNI AaHNn'PHXUheI PXJnatDIeP Xn PhaLePUeaIe AF ZaLXnJ UhN
enSaIJeZentPNC the eSXMaAethan UIXntXnJ that CXJDIed Xn the RNIL. the deUaItZent
NCHXUheIP NC the IXKeIAanL SaANIatNIXeP HNnPXPted NC 14 NI 15 hXJh-PHhNNSand HNSSeJe JIa
RhN aPPXJned the XndXKXdDaS SetteIP Xn thePeeSXMaAethan teYtP tN Nne NI the NtheI NC tRN
\end{lstlisting}

\paragraph{Pasul 4: N=o, C=f.} NC este cuvântul de două litere OF (N a mai apărut în WN = to).
\begin{lstlisting}[language={}]
fIXedZan dXd PoZe JenetXHP RoIL foI hXZ, ADt, AeHaDPe he RaP handFRXth a HaZeIa, he heSU
HIFUtoSoJXPtP Rho ReIe SooLXnJ foI AaHon'PHXUheI PXJnatDIeP Xn PhaLePUeaIe AF ZaLXnJ Uho
enSaIJeZentPof the eSXMaAethan UIXntXnJ that fXJDIed Xn the RoIL. the deUaItZent
ofHXUheIP of the IXKeIAanL SaAoIatoIXeP HonPXPted of 14 oI 15 hXJh-PHhooSand HoSSeJe JIa
Rho aPPXJned the XndXKXdDaS SetteIP Xn thePeeSXMaAethan teYtP to one oI the otheI of tRo
\end{lstlisting}

\paragraph{Pasul 5: X=i.} XG este cuvântul IN (G = n deja cunoscut).
\begin{lstlisting}[language={}]
fIiedZan did PoZe JenetiHP RoIL foI hiZ, ADt, AeHaDPe he RaP handFRith a HaZeIa, he heSU
HIFUtoSoJiPtP Rho ReIe SooLinJ foI AaHon'PHiUheI PiJnatDIeP in PhaLePUeaIe AF ZaLinJ Uho
enSaIJeZentPof the eSiMaAethan UIintinJ that fiJDIed in the RoIL. the deUaItZent
ofHiUheIP of the IiKeIAanL SaAoIatoIieP HonPiPted of 14 oI 15 hiJh-PHhooSand HoSSeJe JIa
Rho aPPiJned the indiKidDaS SetteIP in thePeeSiMaAethan teYtP to one oI the otheI of tRo
\end{lstlisting}

\paragraph{Pasul 6: R=w, P=s.} RTP = WAS, RQN = WHO, RVIV = WERE, după ce T, Q, V sunt cunoscute.
\begin{lstlisting}[language={}]
fIiedZan did soZe JenetiHs woIL foI hiZ, ADt, AeHaDse he was handFwith a HaZeIa, he heSU
HIFUtoSoJists who weIe SooLinJ foI AaHon'sHiUheI siJnatDIes in shaLesUeaIe AF ZaLinJ Uho
enSaIJeZentsof the eSiMaAethan UIintinJ that fiJDIed in the woIL. the deUaItZent
ofHiUheIs of the IiKeIAanL SaAoIatoIies Honsisted of 14 oI 15 hiJh-sHhooSand HoSSeJe JIa
who assiJned the indiKidDaS SetteIs in theseeSiMaAethan teYts to one oI the otheI of two
\end{lstlisting}

\paragraph{Pasul 7: I=r.} CNI = FOR (I este ultima literă a cuvântului).
\begin{lstlisting}[language={}]
friedZan did soZe JenetiHs worL for hiZ, ADt, AeHaDse he was handFwith a HaZera, he heSU
HrFUtoSoJists who were SooLinJ for AaHon'sHiUher siJnatDres in shaLesUeare AF ZaLinJ Uho
enSarJeZentsof the eSiMaAethan UrintinJ that fiJDred in the worL. the deUartZent
ofHiUhers of the riKerAanL SaAoratories Honsisted of 14 or 15 hiJh-sHhooSand HoSSeJe Jra
who assiJned the indiKidDaS Setters in theseeSiMaAethan teYts to one or the other of two
\end{lstlisting}

\paragraph{Pasul 8: J=g.} trigrama XGJ = ING.
\begin{lstlisting}[language={}]
friedZan did soZe genetiHs worL for hiZ, ADt, AeHaDse he was handFwith a HaZera, he heSU
HrFUtoSogists who were SooLing for AaHon'sHiUher signatDres in shaLesUeare AF ZaLing Uho
enSargeZentsof the eSiMaAethan Urinting that figDred in the worL. the deUartZent
ofHiUhers of the riKerAanL SaAoratories Honsisted of 14 or 15 high-sHhooSand HoSSege gra
who assigned the indiKidDaS Setters in theseeSiMaAethan teYts to one or the other of two
\end{lstlisting}

\paragraph{Pasul 9: H=c, F=y, U=p, S=l.} HIFUWNSNJF = CRYPTOLOGY, găsit după șablonul c-r-y-p-t-o-l-o-g-y.
\begin{lstlisting}[language={}]
friedZan did soZe genetics worL for hiZ, ADt, AecaDse he was handywith a caZera, he help
cryptologists who were looLing for Aacon'scipher signatDres in shaLespeare Ay ZaLing pho
enlargeZentsof the eliMaAethan printing that figDred in the worL. the departZent
ofciphers of the riKerAanL laAoratories consisted of 14 or 15 high-schooland college gra
who assigned the indiKidDal letters in theseeliMaAethan teYts to one or the other of two
\end{lstlisting}

\paragraph{Pasul 10: Z=m.} ZTG = MAN.
\begin{lstlisting}[language={}]
friedman did some genetics worL for him, ADt, AecaDse he was handywith a camera, he help
cryptologists who were looLing for Aacon'scipher signatDres in shaLespeare Ay maLing pho
enlargementsof the eliMaAethan printing that figDred in the worL. the department
ofciphers of the riKerAanL laAoratories consisted of 14 or 15 high-schooland college gra
who assigned the indiKidDal letters in theseeliMaAethan teYts to one or the other of two
\end{lstlisting}

\paragraph{Pasul 11: A=b, D=u.} Athng = BACON, ndwsvw = OUTLET.
\begin{lstlisting}[language={}]
friedman did some genetics worL for him, but, because he was handywith a camera, he help
cryptologists who were looLing for bacon'scipher signatures in shaLespeare by maLing pho
enlargementsof the eliMabethan printing that figured in the worL. the department
ofciphers of the riKerbanL laboratories consisted of 14 or 15 high-schooland college gra
who assigned the indiKidual letters in theseeliMabethan teYts to one or the other of two
\end{lstlisting}

\paragraph{Pasul 12: L=k, M=z.} Pqtlvpuvtiv = SHAKESPEARE, Vsxmtavwq = ELIZABETH.
\begin{lstlisting}[language={}]
friedman did some genetics work for him, but, because he was handywith a camera, he help
cryptologists who were looking for bacon'scipher signatures in shakespeare by making pho
enlargementsof the elizabethan printing that figured in the work. the department
ofciphers of the riKerbank laboratories consisted of 14 or 15 high-schooland college gra
who assigned the indiKidual letters in theseelizabethan teYts to one or the other of two
\end{lstlisting}

\paragraph{Pasul 13: K=v, B=q, Y=x, E=j.} Kxijxg = VIRGIN, bdxhl = QUICK, „4 Y 7” = 4 x 7, Enqg = JOHN.
\begin{lstlisting}[language={}]
friedman did some genetics work for him, but, because he was handywith a camera, he help
cryptologists who were looking for bacon'scipher signatures in shakespeare by making pho
enlargementsof the elizabethan printing that figured in the work. the department
ofciphers of the riverbank laboratories consisted of 14 or 15 high-schooland college gra
who assigned the individual letters in theseelizabethan texts to one or the other of two
\end{lstlisting}
